\chapter{Literature Review}
\label{ch:literature-review}

\section{Preliminaries}

Chapter~\ref{ch:introduction} laid out the three limitations of current LLMs that CAEM addresses (stateless operation, poor confidence calibration, absence of self-improvement) and the three architectural commitments (verified episodic memory, confidence-aware routing, constrained self-improvement) that answer them. This chapter supplies the two kinds of background a reader needs to evaluate those commitments in detail. The Preliminaries section that follows establishes foundational concepts and terminology: the transformer backbone, the hallucination taxonomy, natural language inference, sentence embeddings, self-consistency decoding, and catastrophic forgetting. The Review of Existing Research section that follows it then surveys the empirical literature across six areas (hallucination measurement, verification, confidence estimation, memory-augmented architectures, continual learning, and evaluation benchmarks), locating CAEM's contribution relative to prior work. Together they prepare the ground for the formal methodology in Chapter~\ref{ch:methodology}.

\subsection{Large Language Models and Transformer Architecture}

Neural networks that have been trained on quite big text collections to predict and produce natural language are known as large language models. The Transformer architecture \cite{NIPS2017_3f5ee243}, which processes sequences using a self-attention mechanism allowing all the tokens in the input to take account of every other token, is adopted by new LLMs. The encoder and decoder layers in this architecture are arranged in layers and each one has feed forward networks and multi-head self-attention.

In this study, we utilize Qwen-2.5-3B-Instruct \cite{qwen25}, a three-billion-parameter decoder-only transformer that has already undergone instruction-tuning on a broad mixture of task-formatted demonstrations before reaching the CAEM pipeline. Decoder-only models factorise the joint probability of the output tokens left-to-right and produce the answer autoregressively from the question prompt, with the question and the generated rationale co-existing in a single causal attention stream. Qwen-2.5-3B-Instruct is chosen over larger open checkpoints because it is the largest decoder-only backbone that fits the CAEM self-improvement loop on a single consumer-grade GPU when held in sixteen-bit brain-float precision with the MC-Dropout sampling and the nine-signal verifier active on the same device, and it is chosen over smaller decoder-only or encoder-decoder alternatives because its instruction-tuned initialisation makes the cycle-zero verifier outputs sharp enough for the four-outcome storage decision to separate high-purity from borderline items from the first cycle onwards rather than requiring an additional warm-up stage.

Fine-tuning adapts pre-trained models to specific tasks or domains by continuing training on task-specific data with a smaller learning rate. The model parameters $\boldsymbol{\theta}$ are adjusted to reduce loss on the new data while preserving the general abilities learned during pre-training. Standard fine-tuning updates every parameter; parameter-efficient approaches keep the underlying backbone frozen and adapt only a small fraction of the parameter count through trainable side-modules. CAEM ships the parameter-efficient path through the low-rank adaptation layer of \cite{hu2022lora} on the attention and feed-forward projections, with the trainable footprint at roughly two percent of the backbone parameter count and the underlying backbone frozen across every cycle of the self-improvement loop. The empirical evidence motivating that choice over full-parameter fine-tuning is reviewed in \Cref{sec:lit-self-improvement}.

\subsection{Hallucination Taxonomy}
\label{sec:hallucination-taxonomy}

Hallucinations in language model outputs can be understood along two main dimensions. The first dimension separates \textit{intrinsic hallucinations}, which directly contradict information provided in the source input, from \textit{extrinsic hallucinations}, which introduce information that goes beyond what the source provides, regardless of whether it also happens to conflict with the source \cite{ji2023survey}. For instance, if a biography states that ``Einstein was born in Germany'' and the model generates ``Einstein was born in France,'' this is an intrinsic hallucination because it contradicts the stated source. But if the model generates ``Einstein played violin exceptionally well'' without any evidence in the source, this is an extrinsic hallucination even if the claim happens to be true: the defining characteristic is that it cannot be verified from the input, not that it contradicts anything in it.

The second dimension distinguishes \textit{factuality hallucinations}, which conflict with verifiable world knowledge, from \textit{faithfulness hallucinations}, which contradict user input or task context \cite{10.1145/3649506}. If a model claims "photosynthesis produces oxygen from carbon dioxide and water" and demonstrates factual accuracy. But on the other hand it it claims "photosynthesis occurs in mitochondria" then this demonstrates factual hallucination. Faithfulness focuses on whether the output of the model remains consistent with the supplied input or context, especially when summarizing documents or following instructions.

This research focuses on factuality hallucinations, which occur in knowledge-based reasoning tasks. Here, correct answers can be verified against trusted sources or logical rules. Faithfulness to retrieved evidence also matters for retrieval-augmented generation, where outputs must be grounded in retrieved documents.

Within factuality hallucinations, four subclasses are operationally distinguished because CAEM's nine-signal verifier is organised in four families that map onto them. \textit{Confabulations}, in the sense of \cite{farquhar2024semantic}, are outputs the model could equally have produced differently on a resample even when individual token probabilities look confident, and are detected by the sample-set agreement family of two signals ($s_{\text{avg}}$, the mean unique-pair cosine similarity over $M{=}3$ resampled chains in the sentence-embedding space, and $p_{\text{entail}}$, the mean entailment of each resampled chain against the served answer under the natural-language-inference judge).\footnote{The nine signals are defined formally in the Unified Verifier subsection of \Cref{ch:methodology} (table of signal definitions). This chapter refers to them by their canonical symbol only where needed for motivation. The normalised semantic entropy signal of \cite{farquhar2024semantic} was registered earlier in the project alongside the same family but retired before the main run because its cycle-zero per-benchmark Pearson correlation with exact-match labels was indistinguishable from zero on every per-benchmark slice of the calibration fold; the kernel-language-entropy upgrade of \cite{nikitin2024kle} was registered as a candidate replacement and is also retired in the shipped architecture.} \textit{Ungrounded factual claims} are answers without support in retrieved evidence, detected by the external-grounding family of three signals ($p_{\text{ground\_max}}$, $p_{\text{ground\_mean}}$, and the length-gated weakest-link atomic-fact entailment $p_{\text{ground\_atomic}}$ that runs only on hypotheses whose token length exceeds the registered atomic-decomposition gate). \textit{Off-topic answers} that drift from the question are detected by the question-answer relevance signal $q_{\text{a,rel}}$ that scores the served answer against the question through the same cross-encoder used for passage rerank. \textit{Confidently-wrong} factual outputs (internally confident but externally unsupported) are caught by the early-exit confabulation gate that combines a high internal-calibration signal with a low max-match grounding score. The deployed configuration ships MiniCheck~\cite{tang2024minicheck} as the entailment backend across all hypothesis lengths because the long-hypothesis Qwen judge fails the registered Pearson floor on the registered 500-sample calibration overlap fold and is ablated in the locked configuration; the contradiction probability $p_{\text{contra}}$ is structurally zero under MiniCheck (the unary supported-against-context output is the only label class the judge produces), the corresponding hard-veto branch was retired on 2026-04-22, and the contradiction-scoring call was further short-circuited at the verifier compute path so the signal does not enter the gradient or the composite log-odds at all; contradicted factual claims are caught indirectly through the calibrated grounding signals' contribution to the composite. Outside this scope, CAEM does \textit{not} target instruction-following failures, intermediate-reasoning-step errors inside chain-of-thought traces (only the final answer is verified), intrinsic faithfulness hallucinations in summarisation or translation (where the grounding source is the input document rather than world knowledge), or input-contradiction hallucinations that conflict with user-supplied premises. These exclusions set the failure-mode boundary within which the results in Chapter~5 should be interpreted.

\subsection{Natural Language Inference}

Natural Language Inference (NLI) is the task of determining whether a hypothesis sentence logically follows from a premise sentence \cite{10.1007/11736790_9}. Given premise $P$ and hypothesis $H$, an NLI model classifies the relationship as:
\begin{itemize}
    \item \textbf{ENTAILMENT}: $P$ implies $H$ is true
    \item \textbf{CONTRADICTION}: $P$ implies $H$ is false  
    \item \textbf{NEUTRAL}: $P$ neither confirms nor denies $H$
\end{itemize}

For example, if we assume "The Beatles were formed in Liverpool in 1960" and hypothesis is "The Beatles were an English band," then this relationship is ENTAILMENT. Given hypothesis "The Beatles were formed in London," the relationship is CONTRADICTION. But if the given hypothesis is "The Beatles recorded many albums," then the relationship is NEUTRAL.

As the verifier judge in CAEM, we employ MiniCheck-Flan-T5-Large \cite{tang2024minicheck}, a Flan-T5-Large checkpoint fine-tuned on synthetically decomposed language-model-generated claim-support data. MiniCheck is binary supported-versus-unsupported by construction; this is intentional in the deployed configuration, because the Cycle 0 calibration retired the contradiction-veto branch (the hard $p_{\text{contra}}$ rule) once the head-to-head fold confirmed the contradiction signal collapses to zero under the binary judge. A frozen Qwen-3B long-hypothesis judge was registered alongside MiniCheck to handle hypothesis lengths above the MiniCheck encoder's effective range, but it failed the pre-registered Pearson floor on the registered 500-sample calibration overlap fold and is ablated; every hypothesis is therefore dispatched through MiniCheck, with documented 512-token truncation on the rare hypotheses that exceed the encoder range. The legacy generic-NLI baseline RoBERTa-Large-MNLI \cite{liu2019roberta}, fine-tuned on Multi-Genre Natural Language Inference and reaching approximately 90.8\% on the MNLI matched development set, is retained as a Chapter 5 ablation target (the generic-NLI-backend ablation variant registered in \Cref{sec:comp-ablations}) so the backend choice is defended by trajectory evidence rather than by citation alone. The NLI-style judge, whichever backend is active, serves two purposes in CAEM. Firstly, it verifies factual correctness by checking whether generated claims are supported by retrieved evidence. Secondly, it helps group answers that mean the same thing by testing for bidirectional entailment: if $P$ entails $H$ and $H$ entails $P$, then $P$ and $H$ are considered semantically equivalent.

\subsection{Sentence Embeddings and Semantic Similarity}

Sentence embeddings map the length of variable text sequences to fixed dimensional dense vectors in a semantic space where similar meanings have high cosine similarity. Sentence BERT (SBERT) \cite{reimers-gurevych-2019-sentence} produces encasings or embeddings by fine tuning BERT in a siamese network structure on semantic textual similarity tasks. The model we employ, all mpnet base v2, generates 768-dimensional embeddings.

Cosine similarity between vectors $\mathbf{v}_1$ and $\mathbf{v}_2$ is computed as:
\begin{equation}
\text{sim}(\mathbf{v}_1, \mathbf{v}_2) = \frac{\mathbf{v}_1 \cdot \mathbf{v}_2}{\|\mathbf{v}_1\| \|\mathbf{v}_2\|} = \cos(\theta)
\end{equation}
where $\theta$ is the angle between vectors. Similarity ranges from -1 (opposite) to 1 (identical), with 0 indicating orthogonality.

Semantic similarity is different from string similarity. For instance, the questions "Is 17 prime?" and "Is the number 17 a prime number?" have low string overlap but high semantic similarity (approximately 0.98). CAEM uses semantic similarity for key operations like memory retrieval (finding similar questions), novelty detection (avoiding redundant storage), and self-consistency verification (measuring reasoning chain agreement). The system uses FAISS (Facebook AI Similarity Search) \cite{8733051} for fast, efficient approximate nearest-neighbour search, allowing it to retrieve sub-millisecond results over thousands of stored episodes.

\subsection{Self-Consistency Decoding}

Self consistency \cite{wang2023selfconsistency} improves reasoning quality by first generating multiple independent reasoning chains for the same question, then it selects the most common answer. Unlike standard greedy decoding which produces a single output, self-consistency samples $M$ chains using temperature $T > 0$ to introduce diversity.

For CAEM, we adapt self consistency technique for verification rather than choosing an answer. Given question $q$ and generated answer $a$, we generate $M$ additional chains and measure agreement via average pairwise similarity:
\begin{equation}
s_{\text{avg}} = \frac{2}{M(M-1)} \sum_{1 \le i < j \le M} \text{sim}(\mathbf{v}_i, \mathbf{v}_j)
\end{equation}
where $\mathbf{v}_i$ is the embedding of chain $i$, and only unique off-diagonal pairs are used (self-similarity diagonal terms are excluded). High consistency ($s_{\text{avg}} > 0.85$) indicates robust reasoning where multiple paths converge, while low consistency suggests uncertainty or errors. The $s_{\text{avg}}$ symbol is the canonical name for this signal throughout Chapters~3--5; some older literature writes it as $u_{\text{SC}}$. CAEM pairs $s_{\text{avg}}$ with the chain-to-answer entailment signal $p_{\text{entail}}$ (computed on the same $M{=}3$ resampled pool to amortise the decode cost) as the two-signal sample-set agreement family that feeds the calibrated probability composite. The normalised semantic entropy signal of \cite{farquhar2024semantic} was registered earlier in the project as a third sample-set member computed over $K{=}10$ chains at temperature $T{=}1.0$, and the kernel-language-entropy generalisation of \cite{nikitin2024kle} was registered as a candidate replacement that uses graded semantic-similarity kernels rather than hard cluster assignments; both were retired before the main run on the cycle-zero calibration evidence that the entropy signal contributed no discriminative information beyond what $s_{\text{avg}}$ and $p_{\text{entail}}$ on the same resampled pool already provided, so the deployed verifier consumes nine signals across the four families, of which the sample-set agreement family contributes two.

\subsection{Catastrophic Forgetting and Continual Learning}

Catastrophic forgetting \cite{MCCLOSKEY1989109} occurs when neural networks learn new tasks and suddenly lose their ability to perform previously learned tasks. Standard gradient descent focuses only on improving performance for the current task and does not protect earlier knowledge, so parameter updates that help the new task tend to weaken the old ones. Fine-tuning LLMs on domain-specific data can reduce accuracy on general benchmarks by 30--50\% without mitigation strategies.

Elastic Weight Consolidation (EWC) \cite{doi:10.1073/pnas.1611835114} addresses this by identifying essential parameters for previous tasks using the Fisher Information Matrix, then adding quadratic penalties preventing large updates to important parameters:
\begin{equation}
\mathcal{L} = \mathcal{L}_{\text{new}} + \frac{\lambda}{2} \sum_{i} F_i (\theta_i - \theta_i^*)^2
\end{equation}
where $F_i$ measures parameter importance, $\theta^*$ are previous parameters, and $\lambda$ controls regularization strength. A cheaper first-order approximation replaces the full Fisher matrix with an isotropic $L_2$ anchor, $\frac{\lambda}{2}\lVert \boldsymbol{\theta} - \boldsymbol{\theta}_{\text{prev}} \rVert^2$, applied to the adapter parameters of the previous cycle rather than to the full backbone. This penalty is registered in CAEM as the parameter envelope for the inactive full-fine-tune fallback (when it engages, $\lambda > 0$); on the shipped low-rank-adapter path it collapses by design ($\lambda = 0$) because the bounded adapter parameter envelope layered on a frozen backbone is itself the cycle-to-cycle drift bound, and an explicit weight-anchor would otherwise penalise deviation from a zero-anchored full-backbone reference whose weights the loop is not updating. CAEM therefore combines the bounded-adapter envelope with two further mechanisms whose joint effect bounds the cycle-to-cycle drift in the underlying generator. The first is the parameter-efficient adaptation primitive itself: low-rank adaptation \cite{hu2022lora} on the attention and feed-forward projections, with the trainable footprint at roughly two percent of the backbone parameter count and the underlying backbone frozen across the cycle, retains general capability across iterative cycles substantially better than full-parameter fine-tuning at matched training budget on the empirical evidence reviewed in \Cref{sec:lit-self-improvement}. The second is a multi-modal retention guard whose probe set spans a general-knowledge multiple-choice surface, an open-text factoid surface from a held-out training-benchmark test fold, and a commonsense multiple-choice surface from a second held-out training-benchmark test fold; the cycle aborts and rolls the adapter back to the previous snapshot whenever the worst probe ratio against the pristine baseline falls below a fixed floor. The accumulated episodic memory is preserved across rollbacks, so verified knowledge survives an adapter-rollback event even when the weight update does not. This three-mechanism combination addresses the challenge of learning from imperfectly verified data without degrading through error reinforcement.

The conceptual basis for understanding CAEM's nine-signal post-generation verifier, pre-routing confidence, episodic-memory architecture, and constrained self-improvement loop is established by the groundwork above. The following sections describe how prior research develops and validates these individual components, preparing for their integration in Chapter~\ref{ch:methodology}.

\section{Review of Existing Research}

Where the Preliminaries section above defined the concepts and notation (hallucination taxonomy, natural language inference, sentence embeddings, self-consistency, catastrophic forgetting), this section surveys the prior empirical and methodological literature across six areas essential to understanding CAEM's design: hallucination characterization and measurement, verification methods for reasoning validation, confidence estimation techniques, memory-augmented architectures, continual learning approaches, and evaluation benchmarks. Each subsection identifies key contributions and limitations that motivate CAEM's architectural choices, so that a reader who has already absorbed the Preliminaries can move directly to the comparative positioning without re-reading definitional material.

\subsection{Hallucination Problem and Measurement}

Understanding the nature and characteristics of hallucination in LLMs requires methodologies like taxonomic frameworks and quantitative measurement. Early studies on hallucination shows summarization and data to text generation. Ji et al. \cite{ji2023survey} published the first major survey in ACM Computing Surveys. It establishes a study differentiating intrinsic hallucinations that contradicts source material from extrinsic hallucinations. Even though the claims were unverified. Their work shows that hallucinations occurs at several stages. It includes data collection that embeds factual errors, training procedures that reinforce misleading correlations, and inference methods that fail to account for uncertainty.

Huang et al. \cite{10.1145/3649506} extended this study specifically for large language models in ACM Transactions on Information Systems. It is extended by adding the differencs between factuality and faithfulness. Their analysis of causes across the model lifecycle suggests various groundbreaking research. It suggests that effective mitigation should have intervention at several stages instead of relying solely on detection applied after the model has already produced an answer. This directly supports CAEM's multi component strategy which combines memory, verification, and iterative learning.

Quantitative measurement studies of hallucination prevalence reveals concerning statistics across domains and model scales. Lin et al. \cite{lin-etal-2022-truthfulqa} introduced TruthfulQA at ACL2022, a benchmark of 817 questions across 38 categories which has been specifically designed to elicit imitative falsehoods where models reproduce common human misconceptions found in training data. Their result demonstrates an inverse scaling effect: GPT-3 175B achieved only 21-25\% on the truthful and informative metric (or 58\% on truthfulness alone) compared to 94\% human performance. Larger models often performed worse than smaller ones. This proves that only scaling is not sufficient for truthfulness and thus it needs strong motivation for architectural innovation.

The inverse scaling result has major result for hallucination mitigation strategies. If adding parameters reduces truthfulness on some tasks then mitigation strategies must operate beyond parameter growth. CAEM addresses this by incorporating external verification and memory components that operate independently of model scale, providing corrections and verified knowledge regardless of base model size.

Min et al. \cite{min-etal-2023-factscore} introduced FActScore at EMNLP 2023, which evaluates long-form generations by decomposing the answer into atomic claims and reporting the share supported against retrieved sources rather than scoring the answer holistically. The method achieves less than two-percent error against human evaluation, and applied to biography generation reports a fifty-eight percent FActScore for ChatGPT, indicating that even highly capable models produce substantial proportions of unsupported atomic facts. CAEM adopts the FActScore convention at the storage gate: the weakest-link atomic-fact entailment $p_{\text{ground\_atomic}}$ is computed by decomposing the served answer into atomic claims, scoring each against the reranked passages, and propagating the minimum claim-level entailment as the answer's grounding floor. The decomposition is registered with a length gate so that short answers (below the gate's token threshold) bypass the atomic step, on the principle that short answers contain a single claim whose mean-match score is already a faithful summary of grounding; this length-conditioning is what keeps the atomic decomposition's compute cost bounded on the five-benchmark panel.

Domain-specific studies show even higher hallucination rates in specific contexts. Alkaissi and McFarlane \cite{alkaissi2023artificial} asked ChatGPT to respond to 25 clinical questions and found that 69--77\% contained hallucinated medical information. Cheli et al.\ \cite{info:doi/10.2196/53164} studied hallucination in medical-literature queries and reported rates of 69.6\% for GPT-3.5, 28.6\% for GPT-4, and 91.4\% for Google Bard.

Recent benchmark provided comprehensive evaluation across Several dimensions. One of them being legal sector. Factual accuracy has been measured in legal sector. According to Google's FACTS benchmark (2025), it shows that even top performing systems maintain 6.4\% hallucination rates on legal queries. Furthermore, the Vectara Hallucination Leaderboard tracks document summarization hallucination rates ranging from 3\% for GPT-4 to 27.2\% for Google PaLM 2 Chat. It thus demonstrates wide variation among models and confirm that hallucination remains an issue even in state of the art systems.

The information above reveals three critical takeaways that shape CAEM's design. Firstly, hallucination occurs across models, tasks, and domains, with rates ranging from 3\% to over 90\% depending on context. Secondly, larger models do not automatically produce fewer hallucinations, which implies that architectural intervention is necessary. Lastly, verifying information at the level of atomic facts is far more informative than evaluating full responses, motivating CAEM's commitment to fact-level NLI-based verification.

\subsection{Verification Methods}

Verification methods are used to check whether a model's reasoning is correct before the information is stored or shown to users. This subsection explains four approaches (chain-of-thought prompting, self-consistency decoding, semantic entropy, and natural-language inference) that together supply the signals aggregated inside CAEM's composite verifier.

Chain of thought prompting, which was introduced by Wei et al. \cite{wei2022chain} at NeurIPS 2022, improves reasoning quality by prompting models to generate the next steps before final answers. Showing these intermediate steps leads to stronger reasoning and provides better starting points for verification. When tested on the PaLM 540B model, chain of thought greatly improved performance on math reasoning tasks, increasing accuracy on GSM8K from about eighteen percent to over fifty eight percent. The authors found that this technique showed an emergent ability threshold around 100 billion parameters for \textit{base models} trained on raw text, below which chain-of-thought provided minimal benefit. However, this threshold does not apply to instruction-tuned models, which learn to follow step-by-step reasoning patterns during the instruction-tuning phase itself. Chung et al.\ \cite{JMLR:v25:23-0870} demonstrated that Flan-T5 models show substantial chain-of-thought benefits at 780 million parameters after instruction tuning, confirming that the threshold does not apply to instruction-tuned backbones. Since CAEM uses Qwen-2.5-3B-Instruct~\cite{qwen25}, which is a decoder-only instruction-tuned model at comparable scale, chain-of-thought prompting is effective despite being well below the 100B base-model threshold.

For CAEM, chain of thought serves two purposes. Firstly, it improves base model reasoning quality. It also provides better starting points for verification. Secondly, it gives explicit reasoning chains that enables step by step verification through NLI which also checks whether each reasoning step is supported by evidence. The explainability of chain of thought reasoning proves essential for diagnosis when verification fails, enabling error analysis that identifies which reasoning steps caused failures.

Wang et al. \cite{wang2023selfconsistency} expanded this idea in 2023 by introducing self consistency. Instead of producing a single reasoning path, the model generates several different explanations and then selects the most common answer. The key idea is that difficult and complex problems usually has multiple valid reasoning paths which leads to the same correct answer. On the other hand, incorrect answers rarely achieve consistency across independent derivations. Quantitative improvements with PaLM 540B demonstrate substantial gains: +17.9\% on GSM8K, +11.0\% on SVAMP, +12.2\% on AQuA, and +6.4\% on StrategyQA.

CAEM uses self consistency in different ways. It chooses reasoning chain agreement through semantic similarity rather than selecting the majority answer. If the chains mostly agree, the reasoning is then considered reliable and suitable for storage. But if they differin opinion then the system treats the result as uncertain and may apply further checks or reject it. As a result, this method is efficient since it only requires a small number of extra generations which makes it practical for regular verification. And it also provides useful confidence signals.

Another important method is semantic entropy introduced bt Farquhar et al. \cite{farquhar2024semantic} in 2024. It represents the most rigorous treatment of uncertainty based hallucination detection. Their key innovation addresses a primary limitation of token level entropy: Answers that express the same idea in different wording are grouped together using inference checks, and uncertainty is measured across these meaning based groups.

The methodology achieves AUROC of 0.790 for fabrication detection on question answering tasks. It substantially outperforms token entropy (0.691) and verbalized probability methods (0.698). Even though the performance remains robust across model families including LLaMA 2 (7B-70B), Falcon (7B-40B), and Mistral 7B, the Nature publication validates semantic entropy as a cutting edge technique for uncertainty quantification, motivating its inclusion in CAEM despite computational overhead.

Semantic entropy serves as the third verification layer for CAEM. It catches reports where both NLI and self consistency fail. Let us consider a model which is asked "What is the capital of Australia?" that consistently always generates "Sydney" across multiple chains with high individual probabilities. NLI verification may pass if the model generates factually correct supporting statements about Sydney. Here, self consistency shows perfect agreement. But still the semantic entropy will generate 10 diverse samples and will measure cluster entropy. This may reveal that the model also produces "Canberra" and "Melbourne" at high temperature thus indicating underlying uncertainty. This detection capability for systematic confabulation justifies the computational cost.

Natural language inference provides direct verification by checking whether a claim is supported by evidence. Thorne et al. \cite{thorne-etal-2018-fever} introduced the FEVER dataset at NAACL 2018 which formulates fact verification as NLI with 185,445 claims generated from Wikipedia. Systems must classify claim evidence relationships as SUPPORTS, REFUTES, or NOT ENOUGH INFO. This three way classification would map directly to CAEM's storage decisions. The decision would be to support claims merit storage, refute claims are rejected, and insufficient evidence trigger uncertainty handling.

Pre trained NLI models such as RoBERTa-Large-MNLI has achieved approximately 90\% accuracy on the MNLI benchmark. It thus provides reliable verification for factual claims. But the limitation is that even though NLI verifies individual factual statements, they may miss logical reasoning errors that do not contradict facts. For example, a model could state "Sydney is Australia's largest city" which is true and conclude "therefore Sydney is the capital" which is an invalid reference. NLI would verify the factual premise while missing the logical error thus motivating the need for complementary verification layers.

A more recent strand of work brings conformal prediction to the language-model claim level. Mohri and Hashimoto~\cite{mohri2024conformal} introduced \emph{Language Models with Conformal Factuality Guarantees} at ICML 2024, decomposing each generated answer into atomic sub-claims, scoring each, and removing claims below a calibrated threshold $\tau$ such that the surviving sub-claims are jointly correct with probability at least $1-\alpha$ on a small annotated calibration set. Their framework is the closest published method that explicitly crosses an eighty-percent factuality target without filter-time gold labels. Yadkori et al.~\cite{yadkori2024abstention} extended this with conformal-risk-control to bound the hallucination rate at a user-specified target through learned abstention rather than claim removal, providing the direct precedent for CAEM's abstention class in the four-outcome decision tree. Quach et al.~\cite{quach2024conformal} present \emph{Conformal Language Modeling} at ICLR 2024, calibrating a stopping rule, a rejection rule, and a confidence rule simultaneously through Learn-then-Test multiple-hypothesis testing. Cherian et al.~\cite{cherian2024conformal} sharpen the Mohri--Hashimoto framework with a boosted logistic aggregation over multiple per-claim scores, and supply the algorithmic basis for CAEM's boost layer over the per-signal log-odds in \Cref{sec:composite-boost}. CAEM considered the split-conformal formulation as the storage-gate primitive and rejected it on the registered cycle-zero evidence reviewed in \Cref{sec:conformal-gate}: the calibration-versus-evaluation distribution shift on the five-benchmark panel breaks the marginal-coverage premise, with the realised eval-fold precision falling fifteen to fifty percentage points below the conformal target on training benchmarks. The deployed gate is therefore a fixed-threshold rule on the calibrated probability, in the lineage of post-hoc calibration analyses by \cite{niculescu2005predicting} and the selective-classification framing of \cite{geifman2017selective}.

The integration of these methods provides complementary strengths inside CAEM's nine-signal verifier. Chain of thought improves reasoning interpretability and gives the verifier explicit intermediate claims to score. Self-consistency and chain-level entailment contribute the sample-set agreement family of two signals: the average pairwise similarity $s_{\text{avg}}$ over $M{=}3$ resampled chains, and the chain-to-answer entailment probability $p_{\text{entail}}$ averaged over the same resampled pool. NLI against retrieved evidence contributes the external-grounding family of three signals: a max-match entailment probability $p_{\text{ground\_max}}$, a mean-match entailment probability $p_{\text{ground\_mean}}$, and a length-gated weakest-link atomic-fact entailment probability $p_{\text{ground\_atomic}}$ that decomposes the answer into atomic claims under the FActScore convention~\cite{min-etal-2023-factscore} and runs only when the answer's token length exceeds the registered atomic-decomposition gate; below the gate the atomic signal collapses to the mean-match value. The contradiction-probability branch was retired in the locked configuration: MiniCheck's binary supported-versus-unsupported judge yields $p_{\text{contra}} = 0$ structurally, the registered head-to-head calibration confirmed the signal carries no calibrated information once the long-hypothesis Qwen judge is ablated, the corresponding hard-veto rule was removed on 2026-04-22, and the contradiction-scoring call itself was short-circuited at the verifier compute path so the signal does not enter the gradient or the composite. The three external-grounding signals join three internal-calibration signals (token probability, MC-Dropout variance, and a derived internal-state composite), the two sample-set signals, and Branch~C's question-answer relevance signal $q_{\text{a,rel}}$ to give the nine-signal record, whose \emph{calibrated-probability} composite $\ustored$ feeds the four-outcome storage decision through a fixed-threshold rule on the calibrated probability scale. The composite is not a fixed weighted sum: each signal is mapped to a calibrated $\Pr(\text{em}=1\mid\text{signal})$ through per-signal isotonic regression with auto-detected direction~\cite{niculescu2005predicting}, the per-signal curves are blended with a per-benchmark child fit through a James-Stein style shrinkage prior, and the per-signal log-odds are aggregated through an $L_{2}$-regularised logistic boost layer in the lineage of Cherian-Gibbs-Cand\`{e}s conditional boosting~\cite{cherian2024conformal}. The fixed-threshold rule then reads the calibrated composite at two registered cut-points (the storage cut-point and the deferral cut-point) and emits the four outcomes \textsc{Store}, \textsc{Defer}, \textsc{Abstain}, \textsc{Discard}; the realised cut-point values and the architectural cost of the calibration discipline are reported in Chapter~5. The point for this chapter is that the literature supplies four well-validated families of signals whose failure modes are substantially disjoint, which is what justifies a calibrated multi-signal composite over any single-signal detector.

A complementary strand of architectural literature concerns the prompt-design conditions under which a calibrated verifier can score the served answer at all. Two failure modes register as upstream of any signal-level diagnostic. The first, empty-answer non-compliance, occurs when the prompt structure causes the generator to emit empty or whitespace-only completions on a non-trivial share of queries; the verifier then has no hypothesis to score and the calibrated composite degenerates to its prior. The second, open-QA over-abstention, occurs when a refusal-aware prompt template instructs the generator too aggressively to abstain on uncertainty, with the consequence that the storage class is starved of the in-distribution training signal it requires to fit a calibrated composite. CAEM's locked configuration adopts two registered prompt-design fixes against these failure modes: a constrained-emit grammar that requires a non-empty completion in the answer span before the verifier is dispatched, and an open-QA template that elicits a best-effort answer alongside the abstention option rather than instead of it. Both fixes are pre-registered, evaluated against the pre-fix configuration on the same evaluation fold, and the empirical compliance and abstention deltas are reported in Chapter~5.

\subsection{Confidence Estimation}

Reliable confidence estimation helps systems tell the difference between answers that can be trusted and ones that need extra checking from outside sources. This section looks at basic methods for measuring uncertainty and explains how these ideas have been adapted for use with large language models.

Gal and Ghahramani \cite{pmlr-v48-gal16} laid the theoretical foundations for practical uncertainty estimation in their ICML 2016 paper "Dropout as a Bayesian Approximation." Their key contribution explains that dropout usually during training can be understood as approximate Bayesian inference over model parameters. By keeping dropout active during testing and running the model multiple times, it becomes possible to measure how much the outputs vary.This variation can be used to estimate epistemic uncertainty, which refers to uncertainty caused by limited knowledge about the model itself. And this is different from aleatoric uncertainty, which comes from noise or randomness in the data. But it cannot be reduced.

The theoretical framework shows that a network with dropout before each weight layer approximates a deep Gaussian process. This predicted average and variance can be approximated through Monte Carlo sampling. It describes as performing $K$ forward passes with different dropout masks and computing output variance. For better understanding, this provides uncertainty estimates without requiring group training or modifying the loss function. Because this approach works with any network that already uses dropout, it is especially useful for adapting existing large language models.

For CAEM, MC Dropout is one of the three internal-calibration signals inside the post-generation verifier ($u_{\text{token}}$, $u_{\text{dropout}}$, and the derived $u_{\text{internal}} = 0.5\,u_{\text{token}} + 0.5\,(1 - u_{\text{dropout}})$ that combines them). It quantifies model uncertainty through output variance: high variance indicates the model is unstable with respect to its own weights and suggests that the composite score should be pulled down, while low variance indicates confident generation. This signal is insufficient on its own, since models can be confidently wrong, but the computational cost of the registered five forward passes is acceptable inside the verifier, which only runs on queries that have already been dispatched to Tier~2 or Tier~3 generation. MC~Dropout is deliberately \textit{not} used in the pre-routing confidence $\upre$, which must be cheap enough to run on every query before tier selection.

A major problem with neural network confidence is that models are often too confident in their answers. Guo et al. \cite{pmlr-v70-guo17a} demonstrated this at ICML 2017. They showed modern deep networks predict confidence scores substantially which exceeds actual accuracy rates. For example, predictions at 80\% confidence may achieve only 60\% accuracy and create a  dangerous miscalibration for decision making. They introduced Expected Calibration Error (ECE) as the standard metric, dividing predictions into bins by confidence and measuring the gap between confidence and accuracy within each bin.

Their proposed solution, temperature scaling, achieves remarkable effectiveness with minimal complexity. The method divides logits by a learned scalar temperature $T > 0$ optimized on a validation set: $\hat{y} = \sigma(\mathbf{z}/T)$ where $\mathbf{z}$ are logits and $\sigma$ is softmax. Temperature scaling preserves accuracy (argmax unchanged) of classification while calibrating probabilities. Even though only a single parameter is used, it outperforms more complex methods including histogram binning, isotonic regression, and Bayesian approaches. The simplicity and effectiveness make temperature scaling the standard calibration method applied after training is complete.

CAEM uses temperature scaling as the final step in each stage of confidence estimation: once on the pre-routing $\upre$ and once on the post-generation composite $\ustored$. This step adjusts the raw score so that it matches the real accuracy observed on a held-out calibration split, so that when either stage reports eighty percent confidence, about eighty percent of those predictions are actually correct. This makes it possible to set reliable thresholds for tier selection, for the four-outcome storage decision, and for the training-pool filter $\tau_{\text{train}}$.

Kadavath et al. \cite{kadavath2022language} from Anthropic explored self evaluation in large language models at NeurIPS 2022. They introduced two methods for extracting confidence from LLMs themselves. First is P(True) asks models to evaluate the probability their own previous answers are correct. And, P(IK) or "I Know" prompts models to predict confidence before generating answers. Their findings reveal nuanced calibration properties. Larger models tend to show better calibration on multiple choice questions which means their confidence levels line up more closely with how often they are actually correct. However, reinforcement learning from human feedback can make calibration worse, since it often rewards answers that sound helpful or confident rather than ones that are strictly accurate.

The idea that reinforcement learning from human feedback can hurt calibration is important. Models are trained to be helpful. But they may express high confidence to appear authoritative even when uncertain which is exactly the opposite of desired behavior for high stakes applications. Because of this conflict between sounding helpful and being truthful, CAEM relies on several outside confidence signals instead of trusting the model's own confidence alone.

Xiong et al. \cite{xiong2023can} conducted a detailed study of confidence estimation methods for large language models. They provided comprehensive approaches into white box which use internal information like token logits and sequence probabilities, and black box methods, which rely on external signals such as verbalized confidence or consistency checks. Their systematic evaluation across multiple benchmarks reveals striking findings. That is, even the strongest single methods only reached about 0.50 to 0.60 AUROC for telling correct answers from incorrect ones, which is just slightly better than random guessing.

CAEM's multi-signal approach is influenced by this limited discriminative ability of any individual signal. The system splits confidence into two stages. The first stage, pre-routing confidence $\upre$, fuses two cheap signals (a token-probability estimate from a short draft decode, and an encoder-side convergence ratio $C_{\text{conv}}$ over the final attention layers) through a temperature-scaled sigmoid, and is computed on every query. It exists to drive tier selection and the independent safety override. The second stage is the nine-signal post-generation verifier, which runs only after Tier-2 or Tier-3 generation and computes three internal-calibration signals (the mean token log-probability $u_{\text{token}}$, the MC-Dropout variance $u_{\text{dropout}}$~\cite{pmlr-v48-gal16}, and the derived internal-state composite $u_{\text{internal}} = 0.5\,u_{\text{token}} + 0.5\,(1-u_{\text{dropout}})$), two sample-set agreement signals ($s_{\text{avg}}$, $p_{\text{entail}}$ on the same $M{=}3$ resampled pool), three external-grounding signals ($p_{\text{ground\_max}}$, $p_{\text{ground\_mean}}$, $p_{\text{ground\_atomic}}$), and Branch~C's question-answer relevance signal $q_{\text{a,rel}}$ that guards against off-topic generations. These nine signals feed the calibrated-probability composite $\ustored$ through per-signal isotonic regression with a per-benchmark shrinkage prior and the $L_{2}$-regularised boost layer; the contradiction probability $p_{\text{contra}}$ is excluded from the composite signal set in the locked configuration because the MiniCheck-only deployment makes it structurally zero. The realised per-signal isotonic curves, the boost-layer weights, and the fixed-threshold cut-points at the locked configuration are reported in Chapter~5.

Recent work by Tian et al. \cite{tian2023just} questions if models genuinely know when they don't know. It also questions if the models only match patterns for uncertain expressions from training data.Their study found that verbalized confidence often reflects patterns in the training data rather than the model's actual knowledge or uncertainty. This reinforces CAEM's strategy of combining implicit signals like token probability, dropout, and self consistency with explicit methods, instead of relying only on verbalized confidence, which can be misleading or superficial.

A complementary 2024 survey by Geng et al.~\cite{geng2024survey} catalogues white-box and black-box confidence-estimation methods for language models across the same axes that drive CAEM's verifier ensemble, and confirms that no single method exceeds AUROC of approximately $0.6$ across heterogeneous benchmarks; this survey is the foundational positioning document for CAEM's multi-signal design choice. Two recent single-signal methods deserve specific mention as alternatives we considered. Chen et al.~\cite{chen2024inside} introduced INSIDE / EigenScore at ICLR 2024, computing the log-determinant of the covariance matrix over middle-layer hidden states across $K$ resampled responses; the method reports AUROC near $0.84$ on TriviaQA, the strongest unsupervised single-signal result on factual QA we are aware of, and is registered in CAEM as a candidate replacement for the dropout signal that was not adopted in the locked configuration. Nikitin et al.~\cite{nikitin2024kle} introduced Kernel Language Entropy at NeurIPS 2024, generalising Farquhar et al.'s hard-cluster semantic entropy to a soft-kernel von Neumann entropy over pairwise semantic similarities; KLE is registered as a future-work candidate should CAEM re-introduce the semantic-entropy family in a later iteration (the corresponding $h_{\text{norm}}$ signal of \cite{farquhar2024semantic} was retired before the main run on cycle-zero evidence that it contributed no discriminative information beyond $s_{\text{avg}}$ and $p_{\text{entail}}$ on the same resampled pool), but it is not adopted in the locked configuration because the empirical AUROC gain is bounded above the cost of the additional kernel computation. The selective-prediction precision-versus-coverage tradeoff that CAEM reports in \Cref{ch:evaluation} follows the form pioneered by Kamath, Jia, and Liang~\cite{kamath2020selective} for selective question answering under domain shift, where a calibrator over feature signals predicts $P(\text{correct})$ and the selective rule answers only when the calibrator's confidence exceeds the registered precision target.

The confidence estimation literature establishes several key findings for CAEM's design. Firstly, relying on a single confidence signal does not provide enough accuracy for safe routing, which makes combining multiple signals necessary. Secondly, systematic overconfidence requires calibration applied after training through per-signal isotonic regression and through a calibrated-probability composite. Thirdly, white-box methods that access internal model information offer complementary insights compared to black-box approaches based on output consistency. Lastly, RLHF training may degrade calibration, so pre-RLHF models or instruction-tuned checkpoints such as Qwen-2.5-3B-Instruct are expected to give more reliable confidence signals than RLHF-aligned chat models of comparable scale. Together, these points shape CAEM's two-stage confidence architecture. The first stage is a temperature-scaled pre-routing score $\upre$ that fuses token probability with encoder convergence and runs on every query. The second stage is a nine-signal post-generation verifier producing the calibrated-probability composite $\ustored$ through per-signal isotonic regression (with a per-benchmark shrinkage prior toward the pooled fit) and an $L_{2}$-regularised boost layer, whose output is then read by a fixed-threshold rule at two registered cut-points to emit the four-outcome storage decision; both stages are fitted on a held-out calibration split with no live label information at deployment time.

\subsection{Memory-Augmented Architectures}

External memory architectures give neural networks a place to store information outside their usual parameters. It  helps them learn and fetch facts. This subsection explains how these ideas developed. Starting with early systems that used differentiable memory and ending with the modern approach known as retrieval augmented generation.

Graves et al.\ \cite{graves2014neural} introduced Neural Turing Machines (NTMs) in 2014. NTMs ventured into the idea of giving a neural network external memory that it could read from and write to in a smooth, trainable way. The architecture combines a controller network with an addressable memory matrix, supporting both content-based addressing (via cosine similarity) and location-based addressing (via shift operations). The key innovation is a read/write mechanism that enables end-to-end training via backpropagation. As a result, NTMs successfully learned algorithmic tasks including copying, sorting, and associative recall.

The importance of CAEM is that it shows how external memory can give a model abilities that normal models do not have. NTMs were mainly used for learning algorithms. But, it can store facts as well. CAEM uses this idea in its episodic memory. It stores verified question and answer pairs along with the reasoning behind them. This lets the model call them up when needed without any delay.

Weston et al. \cite{weston2015memory} introduced Memory Networks at ICLR 2015. Their model scaled to millions of memory slots. It also showed that a system could perform several rounds of reasoning to answer a question. The architecture consists of four learned components. The components areL: turning the input into features, updating the memory, turning the memory output into features, and finally producing a response. Furthermore, Memory Networks performed strongly on bAbI question. With answering tasks that required reasoning over multiple facts. Sukhbaatar et al. \cite{NIPS2015_8fb21ee7} extended this  model with end to end Memory Networks at NeurIPS 2015. He removed the need for supervision on retrieval. He did this model by continuous attention mechanisms. This described how CAEM builds its own style of step by step reasoning using stored knowledge.

Graves et al. \cite{graves2016differentiable} presented the Differentiable Neural Computer (DNC) in Nature 2016. It built on the idea of Neural Turing Machines but added smarter memory tools like dynamic allocation. As a result it keep track of the order in which things were written. DNCs achieved 3.8\% average error on bAbI tasks compared to what it achieved for memory networks which is 6.4\%. It clearly shows that calculated memory management improves performance.

Lewis et al. \cite{NEURIPS2020_6b493230} introduced Retrieval-Augmented Generation (RAG) at NeurIPS 2020. They showed that combining what a model knows in its parameters with information stored outside the model can make its answers much more accurate. RAG adds a large searchable index of Wikipedia to a BART model with four hundred million parameters. It uses Dense Passage Retrieval to look up relevant documents. Then it works out the final answer by considering the generation probability across all the documents it retrieved. The architecture achieves substantial improvements over parametric-only models: 44.5\% exact match on Natural Questions (versus 41.5\% prior SOTA), 68.0\% on TriviaQA (versus 60.5\%), and 45.5\% on WebQuestions (versus 44.7\%).

Critically, RAG demonstrates index hot-swapping: updating the document index without retraining achieves 70\% accuracy on questions about updated world knowledge, proving that non-parametric memory can be updated without expensive retraining. This property directly motivates CAEM's episodic memory, which can be updated continuously through verification without model fine-tuning, though CAEM additionally employs iterative fine-tuning to consolidate learned patterns.

The limitation of standard RAG for CAEM's purposes is that retrieved knowledge is not verified before use and does not accumulate verified reasoning chains. CAEM extends RAG by adding verification layers and storing verified outputs for future retrieval, creating a growing knowledge base rather than relying solely on static external corpora.

Johnson et al.\ \cite{8733051} developed FAISS to handle similarity search at the scale of billions of items in 2019, enabling efficient implementation of memory-augmented systems. The library achieves an 8.5$\times$ speedup over prior GPU methods through optimised k-selection algorithms. FAISS also supports several index types, such as IVF with product quantisation, which allows the system to search in sub-linear time. For CAEM, FAISS is the tool that makes real-time memory search possible across the low tens of thousands of stored episodes accumulated across the trajectory; keeping search time below one tenth of a second is important for the system to remain usable in practice.

Reimers and Gurevych \cite{reimers-gurevych-2019-sentence} introduced Sentence-BERT (SBERT) at EMNLP 2019. SBERT produces semantically meaningful sentence embeddings through siamese-network fine-tuning. The speed improvement over standard BERT is substantial: finding the most similar sentence pair in a group of ten thousand sentences takes about sixty-five hours with standard BERT, but only about five seconds with SBERT, because one can compare the precomputed embeddings directly instead of running a heavy cross-encoder over every pair. This efficiency gain is what makes CAEM's memory search possible: each query is turned into a 768-dimensional embedding, and the system can look through the entire memory index in real time.

More recent work focuses on writeable episodic stores and on insertion-time memory augmentation, both of which are direct architectural neighbours of CAEM's memory module. Das et al.~\cite{larimar2024das} introduced Larimar at ICML 2024, an associative key-value episodic memory with one-shot edits, selective forgetting, and an order-of-magnitude faster knowledge-editing path than direct fine-tuning; Larimar establishes that a writeable episodic store is a feasible alternative to parametric storage but does not associate a quality scalar with each entry. Xu et al.~\cite{xu2025amem} introduced A-MEM, a Zettelkasten-inspired agentic memory that writes structured textual notes at insertion and evolves links between memories over time; A-MEM is the closest architectural neighbour to CAEM in the sense that both write structured metadata at insertion, but A-MEM's metadata is textual (tags, links) while CAEM's is a single calibrated scalar storage score that drives a tiered router rather than evolving link structure. Jeong et al.~\cite{jeong2024adaptive} introduced Adaptive-RAG at NAACL 2024, a three-tier router that classifies queries by complexity and dispatches to no-retrieval, single-step retrieval, or multi-step retrieval accordingly; Adaptive-RAG shares CAEM's surface tier count and a backbone-family neighbourhood but routes on query-side complexity at inference time, not on insertion-time stored confidence on a shared memory, and that distinction is what CAEM's three-tier router differentiates against in \Cref{sec:router}.

These memory architecture developments establish the feasibility and benefits of external knowledge storage. NTMs prove differentiable memory enables structured reasoning. Memory Networks scale to millions of facts with multi-hop reasoning. RAG demonstrates that non-parametric knowledge improves factual accuracy and supports dynamic updates. FAISS and Sentence-BERT provide the efficient infrastructure for practical deployment. CAEM synthesizes these advances and extends them in three specific ways. First, retrieval is gated by a combined similarity-and-confidence score $\mathcal{S} = \lambda\, s + (1-\lambda)\, \ustored$ rather than similarity alone, so high-similarity matches that were stored with low verifier confidence still route to Tier~3 grounding. Second, routing is overridden by an independent safety gate on the pre-routing confidence $\upre$ that forces Tier~3 whenever $\upre$ falls below the registered safety floor, regardless of how similar the best memory match is. Third, the memory itself is populated only through a four-outcome storage decision (\textsc{Store}, \textsc{Defer}, \textsc{Abstain}, \textsc{Discard}) governed by a fixed-threshold rule on the calibrated probability composite, with a bounded deferred-reconsider buffer (capacity $K_{\text{def}}$, time-to-live $\tau_{\text{ttl}}$) for borderline items, so the knowledge base grows only from items that clear the registered storage cut-point rather than from all retrievals. These extensions distinguish CAEM from static-corpus RAG and motivate the rest of the methodology.

\subsection{Selective and Adaptive Retrieval}
\label{sec:lit-selective-retrieval}

The retrieval-augmented architectures of the previous subsection retrieve unconditionally on every query: every incoming question pays the retrieval cost, and the retrieved passages enter the generator's context whether the query benefits from external evidence or not. A complementary line of work argues that retrieval is not uniformly helpful and proposes architectures that decide whether to retrieve on a per-query basis. The selective-retrieval literature anchors the architectural design of \Cref{sec:ruc}, which adds a binary retrieval-utility classifier between CAEM's three-tier router and its retrieval-augmented fallback so that queries on which retrieval would amplify the question's misconception framing route to the zero-shot tier instead.

Mallen et al.~\cite{mallen2023whennot} established the empirical anchor for the selective-retrieval position at ACL 2023. They introduced PopQA, a fourteen-thousand-question benchmark of entity-centric factoid queries, and showed that large language models handle popular entities through their parametric memory more reliably than through retrieval, while the same models fail systematically on long-tail entities that retrieval can rescue. The architectural lesson is that retrieval utility is a function of the subject entity's popularity rather than a global hyperparameter, and the corresponding routing rule is a per-relation popularity threshold derived from Wikipedia pageview ranks. The adaptive retrieval rule reaches forty-six and a half percent exact match on PopQA against GPT-3 davinci-003, approximately ten absolute percentage points above the non-adaptive baseline that retrieves on every query. The architectural lesson for CAEM is the empirical receipt that retrieval-utility prediction is feasible from cheap question-side features alone (an entity popularity lookup) and does not require a heavy language-model forward pass at routing time, the design constraint that \Cref{sec:ruc} adopts.

Wang et al.~\cite{wang2023skr} extended this line to a learned classifier at EMNLP Findings 2023. Their self-knowledge routing system collects training labels by running each candidate question through both a no-retrieval pipeline and a retrieval-augmented pipeline, recording the per-pipeline exact-match score, and labelling the question as \emph{retrieval-helpful} when the retrieval-augmented exact match exceeds the no-retrieval exact match and as \emph{retrieval-hurt} otherwise. A BERT-based binary classifier is then trained on the labelled set, and at inference the classifier dispatches each query to retrieval or to the zero-shot generator according to its prediction. The training-label rule (compare exact match across the two routings) is the direct precedent for the empirical-delta labelling rule that CAEM's retrieval-utility classifier inherits in \Cref{sec:ruc}; the architectural difference is that CAEM's classifier sits downstream of the existing tier-1 and tier-2 dispatch outcomes and gates only the retrieval-augmented fallback, while Wang's classifier replaces the dispatch entirely.

Jeong et al.~\cite{jeong2024adaptive} introduced Adaptive-RAG at NAACL 2024 as a three-class extension of the binary routing decision. Their classifier (a T5-large language model fine-tuned on automatically-collected complexity labels) dispatches each query to one of three classes (no retrieval, single-step retrieval, multi-step retrieval), trading routing accuracy for the cost of running a large model at routing time. The reported cost-accuracy frontier on a five-dataset open-domain question-answering suite (SQuAD, Natural Questions, TriviaQA, HotpotQA, 2WikiMultiHopQA, MuSiQue) beats single-strategy baselines and the unconditional retrieval-augmented baseline, but the T5-large routing cost is the system's main weakness. The architectural lesson for CAEM is that a learned classifier on automatically-collected outcome labels can decide the routing question, but that putting a large model at routing time costs more than the architecture gains; the follow-up of Moskvoretskii et al.~\cite{moskvoretskii2025llm} extends Adaptive-RAG to the regime where the classifier carries no language-model forward pass and matches the routing accuracy of the language-model-based variant on the same suite. CAEM's retrieval-utility classifier sits at the lightweight-classifier-only end of this design space, with the only language-model dependency being an offline self-knowledge probe that is fit once at cycle zero and consumed as a fixed feature thereafter.

Asai et al.~\cite{asai2024selfrag} introduced Self-RAG at ICLR 2024 as the end-to-end fine-tuned alternative to the external-classifier approach. Self-RAG trains a single language model to emit four kinds of reflection tokens during generation (a \texttt{Retrieve} token that triggers retrieval, an \texttt{IsREL} token that judges passage relevance, an \texttt{IsSUP} token that judges support, and an \texttt{IsUSE} token that judges usefulness), so the routing decision is made by the same model that generates the answer. The deployed signal for retrieval routing is the probability the model assigns to the \texttt{Retrieve} token at decoding time; removing retrieval from the system drops PopQA performance by approximately forty percent relative, the empirical receipt for the joint design's load-bearing role. The architectural cost is that the reflection tokens are tied to the LLaMA-family backbone Self-RAG fine-tunes on, and retargeting them to a different backbone changes enough of the system to invalidate the architectural comparison; this is the same constraint that prevents Self-RAG from entering the external baseline panel of \Cref{ch:evaluation} as a numerical baseline rather than a citation-level reference.

Baek et al.~\cite{baek2025probing} introduced Probing-RAG at NAACL Findings 2025 as a midpoint between the heavy classifier and the in-generator reflection tokens. They train a linear probe on the intermediate hidden states of the generator (rather than on its output tokens), with the probe trained under a balanced positive-negative protocol that explicitly enforces a fifty-fifty class ratio. The probe predicts whether retrieval-augmented generation will improve the answer over zero-shot generation; the reported accuracy across five open-domain question-answering datasets beats FLARE, Self-RAG, and Adaptive-RAG on the same suite with fewer retrieval steps. The architectural lesson for CAEM is the most directly relevant of the five: the closest cousin of CAEM's retrieval-utility classifier on the feature axis is the hidden-state-as-routing-signal pattern that Probing-RAG establishes, with the methodological difference that the latter probes intermediate layers while CAEM consumes the last-token hidden state under the disable-adapter context as a single self-knowledge feature alongside thirty-five other cheap question-side and passage-side features.

The empirical ceiling that this literature collectively establishes for retrieval-utility prediction is a held-out area-under-curve of approximately $0.80$ to $0.86$ in distribution and $0.65$ to $0.75$ out of distribution, with the highest reported in-distribution readings on Probing-RAG and on the SeaKR uncertainty-decoder family. No published system convincingly exceeds an in-distribution area-under-curve of $0.90$ for generalised retrieval-utility prediction on a multi-benchmark held-out fold. CAEM's retrieval-utility classifier targets this ceiling under the lightweight-classifier-only operating point, with a deployed nested-cross-validation area-under-curve of $0.85$ on a three-thousand-two-hundred-and-eight-row held-out fold reported in \Cref{sec:ruc}.

\subsection{Self-Improvement and Continual Learning}
\label{sec:lit-self-improvement}

Self-improvement mechanisms enable models to learn from interaction and feedback, while continual learning addresses the challenge of acquiring new knowledge without forgetting previous capabilities. This subsection reviews methods for iterative refinement and catastrophic forgetting mitigation that inform CAEM's self-improvement loop.

Kirkpatrick et al. \cite{doi:10.1073/pnas.1611835114} introduced Elastic Weight Consolidation (EWC) at PNAS 2017, providing a principled solution to catastrophic forgetting in neural networks. When learning sequential tasks, standard gradient descent optimizes exclusively for current task performance, causing parameter updates that improve new task accuracy while dramatically degrading previous task performance. EWC addresses this by estimating parameter importance for previous tasks using the Fisher Information Matrix, then adding quadratic penalties preventing large updates to important parameters.

The EWC loss function augments standard task loss with a regularization term:
\begin{equation}
\mathcal{L}_{\text{EWC}} = \mathcal{L}_{\text{task}} + \frac{\lambda}{2} \sum_{i} F_i (\theta_i - \theta_i^*)^2
\end{equation}
where $F_i$ approximates the Fisher Information measuring how much loss increases when parameter $\theta_i$ changes, $\theta^*$ are previous task parameters, and $\lambda$ controls regularization strength. Parameters important for previous tasks (high $F_i$) receive strong penalties against change, while less important parameters adapt freely to new tasks.

Applied to Atari game learning, EWC maintained approximately 90\% of single-task performance when sequentially learning 10 games, compared to severe degradation with standard SGD. For CAEM, EWC supplies the conceptual motivation for bounding cycle-to-cycle drift in adapter space; the isotropic $L_{2}$ anchor $\frac{\lambda}{2}\lVert \theta - \theta_{\text{prev}} \rVert^{2}$ that approximates EWC under a uniform-Fisher assumption is registered for the inactive full-fine-tune fallback ($\lambda > 0$), but on the shipped low-rank-adapter path it collapses by design ($\lambda = 0$). The reason is that the bounded adapter parameter envelope layered on a frozen backbone is itself the cycle-to-cycle drift bound: an explicit weight-anchor would otherwise penalise deviation from a zero-anchored full-backbone reference whose weights the loop is not updating, which is meaningless arithmetic. Together with the parameter-efficient adapter envelope and the multi-modal retention guard described below, this bounded-envelope mechanism is sufficient to meet the retention target on the empirical evidence reported in Chapter~5.

The choice to fine-tune through a low-rank adapter rather than the full backbone is a separate and equally load-bearing decision in the self-improvement loop. Hu and colleagues introduced low-rank adaptation \cite{hu2022lora} as a parameter-efficient alternative to full fine-tuning that injects trainable low-rank matrices into the attention and feed-forward projections of a frozen backbone, with the original paper reporting near-full-FT quality at less than one percent of the trainable parameter count on GPT-3 175B. The choice between low-rank adaptation and full fine-tuning has empirical consequences for continual learning that the recent literature has measured directly. Biderman and colleagues sweep Llama-2-7B across rank settings from sixteen to two hundred and fifty-six on Magicoder coding instruction data \cite{biderman2024lora} and report two findings that translate directly to CAEM's setting: low-rank adaptation retains roughly sixty-six percent of HellaSwag, ARC-Challenge, and Winogrande accuracy at every rank tested, against sixty-two percent retention for full fine-tuning matched on epoch count, and the low-rank versus full gap on retention dominates the rank-versus-rank gap on fit. Higher rank trades retention for fit, but every rank setting they tested retained more general capability than full fine-tuning at the same training-pool size. Dettmers and colleagues' QLoRA work \cite{dettmers2023qlora} establishes the second decisive practical point: applying low-rank adaptation to all linear layers, both attention and feed-forward, rather than to the attention projections alone is the most critical configuration choice, and matched-parameter-count experiments show that adding feed-forward adapters to attention adapters dominates adding more attention rank. CAEM's adapter target set spans every linear projection inside the transformer block (the four attention projections together with the three feed-forward projections of the gated multi-layer perceptron), which is the all-linear-layer target set that this literature converges on, fitted at rank thirty-two with the standard $\alpha=2r$ scaling, which sits comfortably inside the rank regime where the standard scaling avoids the gradient collapse at higher rank that motivates the rank-stabilised alternative \cite{kalajdzievski2023rslora}. The low-rank-adaptation primary path is what permits CAEM to fit ten cycles of self-improvement on a single consumer-grade graphics processor in the project's compute envelope; full fine-tuning at the same backbone size would not fit, and the agent-research audit accompanying this thesis registers full fine-tuning as the rejected alternative on the joint retention-and-feasibility axis. The within-adapter analogue of catastrophic forgetting under sequential fine-tuning, documented by Wang and colleagues \cite{wang2023olora} and orthogonal-decomposition fixes for it, is registered as future work; the present thesis combines the bounded low-rank-adapter parameter envelope on a frozen backbone, the multi-modal retention guard with adapter rollback, and the strict training threshold above the storage cut-point as the three-mechanism defence against iterative-adapter drift across the realised cycle trajectory.

The practical implementation in CAEM sets the adapter $L_{2}$ anchor coefficient to $\lambda = 0$ on the shipped low-rank-adapter path (the $\lambda > 0$ branch is registered for the inactive full-fine-tune fallback only) and draws training pairs entirely from verified episodes that exceed a strict training threshold, with a benchmark-balancing reweighting layer that prevents single-benchmark dominance through a temperature-mixed softmax over per-benchmark counts, a registered minimum floor on every benchmark's contribution, a bounded upsampling cap on each benchmark's effective count, and a hard ceiling on any single benchmark's share of the pool. The earlier 90/10 mix that paired episode-specific examples with general-capability examples was retired in favour of the reweighting layer because the adapter's small trainable footprint already bounds drift away from the pretraining mass; the redesign trades a fixed mixing ratio for a benchmark-balanced training pool whose composition adapts to the cycle's actual episode counts. A registered set of held-out probes supplies the retention guard: a general-knowledge probe on the multitask language understanding distribution, a held-out test slice of the open-text factoid-question-answering training benchmark, and a held-out test slice of the multiple-choice commonsense-reasoning training benchmark. Before the first cycle a pristine baseline is recorded for every probe; after each cycle's fine-tune the post-cycle accuracy is recorded and the per-probe ratio against the pristine baseline is computed. The cycle commits the adapter only when the worst probe ratio clears the floor $\rho_{\min} = 0.93$; otherwise the adapter is rolled back to the previous cycle's snapshot and retroactive re-verification is skipped for that cycle. The worst-probe rule replaces the earlier single-probe rule: a single general-knowledge probe missed the failure mode where the multiple-choice accuracy stayed high while open-text generation degraded silently, and the retention guard's purpose is to catch precisely that asymmetric drift before it propagates into the next cycle's gradient. Crucially, the episodic memory is never rolled back even when the adapter update is rejected, so verified knowledge accumulated this cycle remains available for future routing; this asymmetry is deliberate and distinguishes CAEM's retention mechanism from standard continual-learning rollback schemes.

Natarajan et al. \cite{natarajan2013learning} provided theoretical foundations for learning with noisy labels at NeurIPS 2013. Their unbiased estimator method constructs proxy loss functions recovering optimal classifiers despite label corruption. The key insight is that if the noise process is known or can be estimated (e.g., what fraction of positive labels are actually negative), modified loss functions achieve consistent parameter estimation. Their method achieves more than 88\% accuracy with 40\% label corruption which substantially outperforms standard training. 

For CAEM, this framework directly applies to verified episodic memory: no practical verifier is perfect, so a known fraction of stored episodes will carry incorrect labels (false positives that nonetheless pass the calibrated composite). The stricter training threshold $\tau_{\text{train}}$, which is registered strictly above the storage cut-point $\tau_{\text{store}}$ on the calibrated probability scale, is designed exactly to lift the purity of the training pool above the purity of the memory as a whole, and the noise-tolerant-learning literature is what justifies training on the resulting pool at all: it proves that training on pools whose purity is well below 100\% can still improve the model rather than degrading it, provided the corruption rate is below a method-dependent tolerance. The cycle-zero calibration-fold empirical precision at the storage cut-point bounds the pool corruption rate from above in the calibration-consistent sense projected forward by exchangeability between the calibration fold and the cycle's candidate pool, which is the architectural reason CAEM commits to a fixed registered cut-point on the calibrated probability rather than a hand-tuned threshold on raw signals. The realised cut-point, the empirical purity of CAEM's training pool, and the resulting effect on fine-tune outcomes are reported in Chapter~5.

When Zelikman et al. \cite{NEURIPS2022_639a9a17} introduced STaR (Self Taught Reasoner) at NeurIPS 2022, they showed successful bootstrapping using iterative self improvement. STaR loop generates rationales for problems, filters for correct answers, fine tunes on correct rationale-answer pairs and then repeats. A key innovation is rationalization. When the model produces a correct answer without valid reasoning, it generates a rationale conditioned on the correct answer. This data augmentation stops the model from learning problems that it already knows how to solve.

In CommonsenseQA test, STaR achieved +35.9\% improvement over few-shot baseline through iterative refinement. The method shows that bootstrapping from the model's own output works when coupled with appropriate filtering. The direct parallel to CAEM is evident: generate predictions, verify correctness, fine-tune on verified examples, and repeat. CAEM extends STaR in two specific ways. First, simple correctness filtering is replaced by the nine-signal verifier whose calibrated-probability composite $\ustored$ drives the four-outcome storage decision, so the filter is itself a learned verifier rather than a binary match. Second, verified items persist in episodic memory for future retrieval rather than being consumed purely as training data, which turns each cycle into a cumulative knowledge gain rather than a one-off weight update.

The recent work of examining intrinsic self-correction capabilities reveals important limitations. Huang et al. \cite{huang-etal-2023-large} found at EMNLP 2023 that prompting LLMs to review and correct their own responses rarely improves accuracy if we do not use external feedback. Kamoi et al. \cite{kamoi-etal-2024-llms} analyzed when LLMs can correct mistakes at TACL 2024. They found that intrinsic self-correction often fails or degrades performance. The bottleneck is reliable feedback generation. So, models that produce incorrect initial answers usually cannot generate reliable correction signals without external validation. 

This limitation motivates CAEM's architectural approach. Instead of depending on intrinsic self-correction using prompting alone, CAEM employs external verification through an NLI cross-encoder, multi-chain self-consistency, and semantic-entropy analysis over stochastic re-samples. These external signals provide the reliable feedback the self-correction literature identifies as a precondition for effective self-improvement, and their failure modes are sufficiently disjoint that aggregating them into a single composite score is expected to be substantially more reliable than any one of them alone, and substantially more reliable than the model's own self-evaluation. The end-to-end verifier accuracy is reported in Chapter~5. Together with the four-outcome storage decision and the training threshold $\tau_{\text{train}} > \tau_{\text{store}}$, this enables the virtuous improvement cycle in which verified generation produces higher-purity training data, which in turn improves generation.

The self-improvement literature also points out successful cases. When external feedback is available (for example code execution results, formal verification, retrieval of sources), models can effectively adapt the corrections. CAEM's Tier 3 routing with RAG provides exactly the same external feedback by retrieving evidence that grounds generation and enables meaningful verification. The system learns not just from its own outputs but from the interaction between generation and evidence based verification.

Domain adaptation research provides additional relevant findings. Gururangan et al. \cite{gururangan-etal-2020-dont} demonstrated at ACL 2020 that continued pre-training on domain-specific unlabeled data (domain-adaptive pretraining) followed by task-specific fine-tuning outperforms direct fine-tuning on limited labeled data. While CAEM does not perform domain-adaptive pretraining due to computational constraints, the finding supports iterative refinement: progressive training on accumulated verified examples enables gradual domain adaptation.

The distinction between CAEM's approach and standard continual learning deserves emphasis. Traditional continual learning addresses learning sequential tasks A, B, C without forgetting previous tasks. CAEM's scenario differs: the task remains constant (question answering), but the model progressively improves at that task through accumulated verified examples. This setting is closer to curriculum learning and iterative refinement than task-sequential continual learning, though catastrophic forgetting concerns still apply since fine-tuning on domain-specific verified examples risks degrading general capabilities.

Several 2024--2025 contributions tighten the theoretical and architectural framing of self-improvement in ways directly relevant to CAEM's theorems and to its retention guard. Hosseini et al.~\cite{hosseini2024vstar} introduced V-STaR at COLM 2024, training a DPO verifier over correct and incorrect self-generated solutions and using it to rank test-time candidates; CAEM generalises V-STaR by using the verifier for training-data admission rather than just test-time ranking, and by adding the $L_{2}$ anchor and the episodic memory. Song et al.~\cite{song2024gap} formalise the \emph{generation--verification gap} and prove empirically that iterative self-improvement succeeds if and only if the verifier is more accurate than the generator on the noisy-label-recovery task; this is the conceptual basis for CAEM's $\alpha > 1/2$ verifier-purity precondition in \Cref{sec:theory}. Fu et al.~\cite{fu2025collapse} prove theoretically that a verifier filter over synthetic data prevents model collapse in self-consuming training loops and admits both near-term improvement and long-term convergence guarantees; this is the strongest published safety anchor for CAEM's storage-gate-filtered SIL. Das et al.~\cite{das2025retraining} prove at ICML 2025 that retraining on verifier-correct hard labels strictly increases population accuracy in the binary linearly-separable classification setting; this is the closest published analogue of CAEM's pool-purity inequality (\Cref{thm:purity}), and CAEM extends it from binary classification with internal self-prediction to open-ended generation with an external multi-signal verifier. Huang, Block, Foster et al.~\cite{huang2025sharpening} formalise self-improvement as a sharpening mechanism at ICLR 2025 and prove convergence-style amortisation-via-self-training bounds that are the closest published analogue of CAEM's bounded-band stationarity result on the committed-cycle subsequence (\Cref{thm:convergence}). Lang, Vijayaraghavan, and Sontag~\cite{lang2024w2s} provide a complementary expansion-based bound at NeurIPS 2024 for when a strong student trained on weak-teacher pseudolabels outperforms the teacher; their bound requires expansion assumptions rather than a verifier-accuracy threshold and is therefore complementary to, rather than redundant with, CAEM's data-purity formulation.

Synthesis of these methods informs CAEM's self-improvement loop design. EWC supplies the conceptual motivation for bounding cycle-to-cycle drift in adapter space; the explicit $L_{2}$ anchor that approximates EWC under a uniform-Fisher assumption is registered for the inactive full-fine-tune fallback only, and on the shipped low-rank-adapter path the bounded adapter parameter envelope on a frozen backbone acts as the implicit drift bound without the hand-tuned anchor coefficient that the full-parameter fine-tune would require. The low-rank-adaptation primitive bounds the trainable footprint at roughly two percent of the backbone parameter count and frozen-backbone semantics, which the parameter-efficient-fine-tuning literature shows retains general capability across iterative cycles substantially better than full fine-tuning at matched training budget. Noisy-label learning provides theoretical justification for training on imperfectly verified data and motivates the stricter training threshold $\tau_{\text{train}} > \tau_{\text{store}}$ that boosts training-pool purity relative to memory purity. STaR demonstrates successful bootstrapping patterns and motivates the cycle structure. The self-correction research validates the need for external verification rather than intrinsic review and motivates the nine-signal verifier over any single-signal confidence proxy. Finally, the multi-modal retention guard with worst-probe floor $\rho_{\min} = 0.93$ and adapter rollback (but not memory rollback) closes the loop by converting the catastrophic-forgetting risk into a detectable, rejectable event rather than a silent degradation, with a probe set spanning a general-knowledge multiple-choice surface, an open-text factoid surface, and a commonsense multiple-choice surface so that asymmetric drift across answer formats is caught at the cycle boundary. The resulting improvement trajectory across the realised cycle horizon, and the retention margin actually achieved across the multi-probe panel, are reported in Chapter~5.

\Cref{tab:positioning} positions CAEM against the closest published systems on the architectural axes that distinguish them: whether the system stores a per-entry confidence at insertion, what signal the router uses, what filter the training-data admission applies, what regulariser preserves general capability under self-improvement, and whether a formal precision or convergence guarantee is supplied. No single prior system shares CAEM on every axis simultaneously; the table makes the differentiation explicit.

\begin{table}[!htbp]
\centering
\footnotesize
\begin{tabular}{@{}p{2.4cm}p{1.7cm}p{2.6cm}p{2.5cm}p{2.0cm}p{2.0cm}@{}}
\toprule
\textbf{System} & \textbf{Stored conf.} & \textbf{Routing signal} & \textbf{Training filter} & \textbf{CL guard} & \textbf{Formal guarantee} \\
\midrule
\textbf{CAEM} (this work) & yes (calibrated $\ustored$) & combined sim $+$ $\ustored$, $\upre$ override & 9-signal cal-prob composite $+$ fixed-threshold gate (per-bench shrinkage; LoRA SIL) & bounded LoRA adapter envelope $+$ multi-modal retention rollback & calibration-consistent precision floor; pool-purity, monotonicity, convergence, asymptotic-floor theorems \\
A-MEM~\cite{xu2025amem} & no (textual tags) & similarity $+$ link evolution & --- & --- & --- \\
Larimar~\cite{larimar2024das} & no (key-value addr.) & key match & one-shot edit & --- & --- \\
Adaptive-RAG~\cite{jeong2024adaptive} & no & query complexity classifier & --- & --- & --- \\
V-STaR~\cite{hosseini2024vstar} & no & --- & DPO verifier (test-time rank) & --- & --- \\
Mohri-Hashimoto~\cite{mohri2024conformal} & no & --- & split-conformal sub-claim filter & --- & conformal factuality \\
MiniCheck~\cite{tang2024minicheck} & no & --- & NLI judge (fact verification) & --- & --- \\
\bottomrule
\end{tabular}
\caption[Positioning of CAEM against the closest prior systems.]{Positioning of CAEM against the closest prior systems. Insertion-time stored confidence, combined-score routing on a shared memory, the calibrated-probability composite at a fixed-threshold gate, the LoRA self-improvement primitive with adapter-rollback continual-learning guard against a multi-modal retention probe panel, and the joint calibration-consistent precision floor and convergence guarantees together constitute CAEM's defensible novelty axis.}
\label{tab:positioning}
\end{table}

\subsection{External Baselines and Competitive Positioning}
\label{sec:external-baselines}

The related work reviewed above clusters into three families relevant to CAEM's evaluation: prompt-level reasoning methods, retrieval-augmented generation methods, and training-time self-improvement methods. Chapter~5 instantiates one representative from each family as an external baseline so that every mechanism CAEM adds on top of a plain language model is compared against at least one prior system that already uses that mechanism in isolation. This subsection establishes what each baseline is, what it isolates, and what a fair comparison would look like; the numerical results appear in Chapter~5.

The prompt-level family contributes three baselines. The first, \textbf{zero-shot Qwen-2.5-3B-Instruct}, uses the same decoder-only backbone CAEM fine-tunes, with no CoT prefix, no retrieval, and no verification. This is the empirical floor: any gain CAEM reports must exceed the zero-shot score on the same backbone, otherwise the gain is attributable to model capacity rather than to the CAEM mechanism. The second, \textbf{Chain-of-Thought} prompting~\cite{wei2022chain}, adds the standard ``Let's think step by step'' trigger without any further modification. CoT isolates the contribution of intermediate reasoning on the same backbone and controls for the possibility that CAEM's gains reflect better rationale generation rather than the verifier-memory architecture. The third, \textbf{Five-shot Chain-of-Thought}, prepends five in-context $(\text{question}, \text{reasoning}, \text{answer})$ demonstrations from the TriviaQA train split before the CoT trigger and is the strongest prompting-only comparison against CAEM's retrieval-free Tier-2.

The retrieval family contributes three baselines. \textbf{DPR-RAG}~\cite{karpukhin-etal-2020-dense, NEURIPS2020_6b493230} retrieves the top $k=5$ Wikipedia passages from a dense DPR index, concatenates them into a 384-token context, and generates from the same Qwen-2.5-3B-Instruct backbone through the Tier-3 prompt builder; this isolates the contribution of external grounding without any verification or memory. \textbf{CoT$+$DPR-RAG} combines the CoT trigger with DPR retrieval and isolates the contribution of reasoning over retrieved evidence. \textbf{FLARE}~\cite{jiang-etal-2023-active} extends DPR-RAG with active retrieval: rather than retrieving once at the start of generation, FLARE generates one sentence at a time and triggers a retrieval whenever the token-level confidence on the upcoming sentence falls below a registered threshold, with retrieved passages prepended to the generation context for the next sentence; this isolates the contribution of selective and step-conditioned retrieval over single-pass retrieval. A fourth retrieval system, \textbf{Self-RAG}~\cite{asai2024selfrag}, is acknowledged as the most ambitious retrieve-and-critique pipeline in the literature but is included only as a citation-level reference rather than as a replicated numerical baseline, because Self-RAG's critic tokens are tied to a LLaMA-family backbone and retargeting them to Qwen-2.5-3B-Instruct would change enough of the system to invalidate the architectural comparison.

The training-time family contributes one baseline. \textbf{Vanilla fine-tuning} trains Qwen-2.5-3B-Instruct on the verified episode pool for the same realised cycle horizon as CAEM at matched per-benchmark chunk size, but without the multi-modal retention probe panel and rollback guard, without the bounded low-rank-adapter envelope, and without the episodic memory or the verifier at inference time, so every selected episode is treated as a clean label. This isolates the contribution of CAEM's retention-protection machinery: if vanilla fine-tuning degrades on the retention probe panel or on the later benchmark cycles while CAEM does not, the retention guard and the bounded adapter envelope are responsible. The confidence-calibration family contributes one baseline. \textbf{Semantic entropy}~\cite{farquhar2024semantic} computes a model-internal uncertainty score by sampling multiple generations under temperature and clustering them into semantically-equivalent groups, then reports the entropy of the cluster distribution as the per-query uncertainty signal; the baseline runs the same backbone as the zero-shot system and reports the same answer, but exposes an uncertainty surface that is the most direct external comparator to CAEM's multi-signal calibrated probability composite. \textbf{STaR}~\cite{NEURIPS2022_639a9a17} is discussed above as the closest prior method in spirit and is run only as an optional ceiling reference conditional on remaining compute budget: its rationalisation step and its binary correctness filter differ enough from CAEM's nine-signal verifier and four-outcome decision that a large gap to STaR would mix architectural contribution with rationalisation-quality artefacts, so STaR is reported where possible as an aspirational ceiling for what a pure iterative-refinement scheme can reach rather than as a primary competitor.

Taken together, these eight running external baselines (B1--B8 in \Cref{sec:comp-baselines}) span the space of prior systems that share at least one mechanism with CAEM: three isolate the prompt-level component (B1 zero-shot, B2 CoT, B5 five-shot CoT), three isolate the retrieval component (B3 DPR-RAG, B4 CoT$+$DPR-RAG, B7 FLARE), one isolates the training-time component (B6 vanilla fine-tuning), and one isolates a sampling-based uncertainty calibration on the same backbone (B8 semantic entropy). Self-RAG is referenced at citation level only, not run as a numerical baseline, because retargeting its critic tokens to Qwen-2.5-3B-Instruct would change enough of the system to invalidate the architectural comparison; STaR is reported separately as an optional ceiling when compute permits, and neither system enters the significance panel's Holm correction. No single baseline combines all four families, which is precisely the gap CAEM occupies and which Chapter~5's retrieval-utility-classifier on/off ablation isolates against the remainder of the architecture.

\subsection{Evaluation Benchmarks}

Rigorous evaluation requires benchmarks testing diverse reasoning capabilities with objective correctness criteria. This subsection reviews the five primary factual-QA benchmarks that constitute CAEM's locked evaluation panel (FEVER, TriviaQA, CommonsenseQA training, and TruthfulQA, StrategyQA transfer) and the MMLU general-knowledge probe that sits inside the multi-modal retention guard, explaining how each tests a different aspect of hallucination reduction. Two benchmarks registered earlier in the project (Natural Questions and ARC-Challenge) are described in this section for completeness, alongside HotpotQA, but were retired before the main run on the cycle-zero diagnostic evidence that the verifier-balanced-accuracy precondition was violated; all three are kept on disk as registered exclusion evidence rather than carried forward into the live panel.

Thorne et al.~\cite{thorne-etal-2018-fever} created the FEVER dataset at NAACL 2018 for fact verification, providing 185{,}445 claims generated from Wikipedia requiring classification as SUPPORTS, REFUTES, or NOT ENOUGH INFO relative to Wikipedia evidence. Claims are specifically designed to test factual knowledge, such as ``Barack Obama was the 44th President of the United States'' (SUPPORTS) or ``The Beatles formed in Manchester'' (REFUTES, they formed in Liverpool). The NOT ENOUGH INFO category is particularly valuable, testing whether systems recognise insufficient evidence rather than hallucinating conclusions. FEVER directly evaluates the external-grounding family of CAEM's verifier, because the task structure matches the premise-hypothesis entailment check that produces $p_{\text{ground}}$ and $p_{\text{entail}}$. FEVER is also the benchmark on which the four-outcome storage decision is most directly exercised: SUPPORTS claims map to \textsc{Store}, REFUTES claims to \textsc{Discard} via the low composite arising from contradicted retrieved evidence, and NOT ENOUGH INFO claims to \textsc{Abstain} via the low-composite-and-low-grounding rule.

Joshi et al.~\cite{joshi-etal-2017-triviaqa} introduced TriviaQA at ACL 2017, providing 650{,}000 trivia question-answer pairs drawn from quizzing leagues, each paired with multiple supporting documents from the web and Wikipedia. Questions span an unusually wide topical range (history, science, geography, entertainment) and the provided evidence is noisy (only a subset of the retrieved documents actually answers the question), making TriviaQA a realistic test of retrieval-augmented generation under distractor documents. For CAEM, TriviaQA tests whether the Tier-3 retrieval-augmented path can ground answers against noisy evidence and whether the mean-match grounding signal $p_{\text{ground\_mean}}$ remains calibrated when retrieval quality varies across queries.

Kwiatkowski et al.~\cite{kwiatkowski-etal-2019-natural} introduced Natural Questions (NQ) at TACL 2019, providing real anonymised queries to the Google search engine paired with Wikipedia-page long answers and span-level short answers. Unlike synthetic benchmarks, NQ captures the distribution of questions users actually ask, with the accompanying variability in specificity, ambiguity, and presupposition. NQ is the dominant open-domain QA benchmark in recent retrieval-augmented-generation literature~\cite{NEURIPS2020_6b493230, karpukhin-etal-2020-dense}. CAEM registered NQ as a candidate panel benchmark and then retired it before the main run on the cycle-zero diagnostic evidence that the verifier-balanced-accuracy precondition for self-improvement was violated on bare-entity short-answer questions; NQ is kept on disk as registered exclusion evidence rather than carried forward into the live five-benchmark panel.

Lin et al.~\cite{lin-etal-2022-truthfulqa} designed TruthfulQA specifically to measure how models mimic human falsehoods, providing 817 questions across 38 categories in health, law, finance, politics, and conspiracies. Questions are adversarially constructed to elicit common misconceptions: ``What happens if you crack your knuckles?'' where the false but common answer is ``you get arthritis'' and the true answer is ``nothing harmful happens; no proven link to arthritis.'' The inverse scaling finding (larger models often less truthful) makes TruthfulQA the most challenging benchmark for parametric scaling approaches. Human performance reaches 94\% while GPT-3 175B achieves only 21--25\% on the truthful-and-informative metric. For CAEM, TruthfulQA is the benchmark where the sample-set family of signals is expected to pay off: misconceptions typically produce high semantic entropy because the model generates both the correct answer and the misconception with similar confidence across stochastic decodes, and the deployed sample-set agreement signals $s_{\text{avg}}$ and $p_{\text{entail}}$, the surviving proxy for the semantic-entropy principle~\cite{farquhar2024semantic}, are designed to catch this.

Geva et al.~\cite{10.1162/tacl_a_00370} introduced StrategyQA at TACL 2021, presenting 2{,}780 questions requiring implicit reasoning where decomposition steps are not mentioned in questions. For example, ``Could a llama birth twice during the War in Vietnam?'' requires knowing Vietnam War duration (implicit: 1955--1975) and llama gestation period (implicit: approximately 11 months), then computing whether two births fit the timeframe. Human accuracy reaches 87\% while best models achieve approximately 66\%. The implicit reasoning requirement tests whether models genuinely understand underlying concepts or merely perform shallow pattern matching. For CAEM, StrategyQA evaluates the sample-set agreement signal $s_{\text{avg}}$ because implicit reasoning should converge across multiple chains if the underlying world-knowledge is consistent; divergence across chains is a strong hallucination signal specifically on this benchmark.

Clark et al.~\cite{clark2018think} introduced ARC-Challenge at AI2 in 2018, a subset of 2{,}590 grade-school science questions filtered to be hard for co-occurrence and information-retrieval baselines. Questions require multi-step reasoning over science concepts and cannot be answered by surface-level lookup. CAEM registered ARC-Challenge as a candidate transfer benchmark and then retired it before the main run on cycle-zero diagnostic evidence that scientific multi-step reasoning at the three-billion-parameter scale produced verifier-balanced-accuracy below the registered floor; ARC-Challenge is kept on disk as registered exclusion evidence and is not part of the live five-benchmark panel.

Hendrycks et al.~\cite{hendrycks2021measuring} introduced MMLU in 2021, a 57-subject multiple-choice benchmark covering humanities, STEM, social sciences, and professional domains. MMLU is used in CAEM as one of three components of the \textit{multi-modal retention probe panel}, not as a target benchmark. The fine-tune is not optimised for MMLU; instead, MMLU is probed before the first cycle (the pristine baseline) and after every cycle's fine-tune, alongside two further probes drawn from held-out test slices of the training panel: an open-text factoid probe on the held-out TriviaQA test fold, and a commonsense multiple-choice probe on the held-out CommonsenseQA test fold. The cycle aborts and rolls the adapter back to the previous snapshot whenever the worst per-probe ratio against the pristine baseline falls below the fixed floor $\rho_{\min} = 0.93$; the worst-probe rule replaces an earlier single-probe rule because a general-knowledge multiple-choice probe alone missed the failure mode where MCQ accuracy stayed high while open-text generation degraded silently. The episodic memory is preserved across rollbacks. The choice of MMLU as one of the three probes is deliberate: catastrophic forgetting manifests most visibly as erosion of capability on tasks the model was not trained on this cycle, so a broad out-of-distribution probe is a more sensitive trigger than in-distribution QA, and the held-out training-bench test folds add the asymmetric-drift sensitivity that the broad probe alone would miss.

Benchmark complementarity ensures comprehensive evaluation across the registered five-benchmark panel. FEVER tests entailment-based grounding, including the mapping from NOT ENOUGH INFO to the \textsc{Abstain} outcome. TriviaQA tests retrieval-augmented generation under noisy evidence. CommonsenseQA tests commonsense multi-step reasoning over short multiple-choice options. TruthfulQA exercises the semantic-entropy proxy and the q-a relevance signal on imitative falsehoods, as a transfer-only benchmark held out from the SIL training pool. StrategyQA exercises sample-set agreement on implicit multi-hop reasoning, also as a transfer-only benchmark. MMLU, the open-text TriviaQA test fold, and the multiple-choice CommonsenseQA test fold close the catastrophic-forgetting loop through the multi-modal retention probe panel. Together, they cover the spectrum of factual reasoning tasks where hallucination reduction matters most, and success across the panel demonstrates that CAEM's improvements generalise across reasoning types rather than overfitting to specific task characteristics. HotpotQA, Natural Questions, and ARC-Challenge were registered earlier in the project but retired before the main run on cycle-zero diagnostic evidence that the verifier-balanced-accuracy precondition for self-improvement was violated; all three are kept on disk as registered exclusion evidence.

\section{Summary of Key Findings}

The literature review both motivates and supplies the technical premises for CAEM's architectural orientation. Hallucination-measurement studies report prevalence rates varying widely by domain, reaching above 90\% in worst-case clinical settings \cite{alkaissi2023artificial}. TruthfulQA's inverse-scaling finding \cite{lin-etal-2022-truthfulqa}, in which GPT-3 175B achieved only 21--25\% on the combined truthful-and-informative metric, is the clearest demonstration that parametric scaling alone is insufficient for truthfulness, and FActScore \cite{min-etal-2023-factscore} shows that even competent models support only about 58\% of their atomic facts, which is why CAEM verifies atomic facts rather than holistic responses. Together these findings motivate a multi-component architectural intervention rather than a pure scaling approach.

The verification literature supplies four families of signals with substantially disjoint failure modes. Chain-of-thought prompting \cite{wei2022chain} produces the intermediate claims that downstream verification operates on. Self-consistency \cite{wang2023selfconsistency} provides the sample-set agreement signal $s_{\text{avg}}$. Chain-to-answer entailment under the natural-language-inference judge provides the second sample-set signal $p_{\text{entail}}$, computed on the same $M{=}3$ resampled pool to amortise the decode cost. The normalised semantic entropy of \cite{farquhar2024semantic, kuhn2023semantic} and the kernel-language-entropy generalisation of \cite{nikitin2024kle} were registered earlier in the project as candidate sample-set members but retired before the main run on cycle-zero evidence that they contributed no discriminative information beyond what $s_{\text{avg}}$ and $p_{\text{entail}}$ already provide on the same resampled pool. Claim-support verification against retrieved evidence provides the external-grounding family of three signals ($p_{\text{ground\_max}}$, $p_{\text{ground\_mean}}$, length-gated atomic-fact $p_{\text{ground\_atomic}}$). The thesis uses MiniCheck-Flan-T5-Large \cite{tang2024minicheck} as the deployed entailment backend across all hypothesis lengths because the registered long-hypothesis Qwen judge fails the pre-registered Pearson floor on the registered 500-sample calibration overlap fold and is ablated in the locked configuration; MiniCheck is fine-tuned specifically on language-model-generated claim-support data and produces well-calibrated supported-versus-unsupported probabilities on the distribution of claims the verifier actually scores, whereas generic NLI models trained on MultiNLI and SNLI (including RoBERTa-Large-MNLI~\cite{liu2019roberta}) exhibit systematic miscalibration on language-model hallucinations: \cite{semanticillusion2025} report a one-hundred-percent false-positive rate at ninety-five-percent recall for a near-state-of-the-art generic NLI model on HaluEval~\cite{li-etal-2023-halueval}, corresponding to a verifier balanced accuracy at or below the chance threshold, which would falsify the purity theorem's empirical premise. The legacy backend is retained in Chapter~5's verifier-calibration diagnostic so the backend choice is defended by trajectory evidence rather than citation alone. CAEM organises nine signals into a post-generation verifier across four families (three internal-calibration signals, two sample-set agreement signals, three external-grounding signals, and Branch~C's question-answer relevance signal $q_{\text{a,rel}}$), aggregates them through per-signal isotonic regression with a per-benchmark shrinkage prior toward the pooled fit and an $L_{2}$-regularised boost layer into the calibrated-probability composite $\ustored$, and reads the composite through a fixed-threshold rule at two registered cut-points to emit the four-outcome storage decision (\textsc{Store}, \textsc{Defer}, \textsc{Abstain}, \textsc{Discard}). A confabulation early-exit gate in which the sharply-asymmetric $(u_{\text{internal}}, p_{\text{ground,max}})$ conjunction is checked before the composite catches the confidently-wrong failure mode that would otherwise route through the gate. The contradiction probability $p_{\text{contra}}$ is excluded from the composite signal set in the locked configuration because the MiniCheck-only deployment makes it structurally zero, the corresponding hard-veto rule was retired on 2026-04-22, and the contradiction-scoring call itself was short-circuited at the verifier compute path so the signal does not enter the gradient or the composite log-odds.

Composite verification approaches that fuse multiple individual signals are recent but published precedent exists. \cite{valentin2024hallucination} publish the closest analogue: an eight-signal isotonic-stacked hallucination detector on TriviaQA, FEVER, HaluEval, and BIG-bench, with a Spearman cross-signal correlation heat-map that identifies three redundancy clusters among their signals. \cite{vashurin2024polygraph} benchmark twenty-eight uncertainty-estimation methods head-to-head and report that sample-diversity methods such as semantic entropy dominate open-generation tasks such as TriviaQA and CoQA, while verbalised-uncertainty and maximum-softmax-probability methods dominate constrained-output tasks such as MMLU. This heterogeneity across output regimes is the principled justification for composite verification over single-signal detection, and motivates the five-benchmark panel evaluated in Chapter 5 rather than a single-benchmark target. The nine-by-five Spearman matrix across CAEM's five benchmarks (scheduled in Chapter 5 as a Phase 1 Full deliverable) is a direct extension of the Valentin 2024 four-benchmark-by-eight-signal matrix to the retrieval and MC-Dropout signal families that their registry did not cover.

The confidence-estimation literature shows that single signals only reach roughly 0.5--0.6 AUROC for discriminating correct from incorrect answers \cite{xiong2023can}, that modern networks are systematically overconfident and require calibration applied after training \cite{pmlr-v70-guo17a}, and that RLHF can further degrade calibration \cite{kadavath2022language}, which is one reason CAEM builds on Qwen-2.5-3B-Instruct's plain instruction-tuned checkpoint rather than an RLHF-aligned chat variant. These findings shape CAEM's two-stage confidence architecture: a cheap temperature-scaled pre-routing score $\upre$ that fuses token probability with encoder convergence and runs on every query, and the nine-signal post-generation verifier whose calibrated-probability composite $\ustored$ runs only after generation. The pre-routing score drives tier selection through the combined score $\mathcal{S} = \lambda\,s + (1-\lambda)\,\ustored$ and an independent safety override that fires whenever $\upre$ falls below the registered safety floor $\phi_{\text{safe}}$.

Memory-augmented architectures establish the feasibility and benefit of external knowledge storage. Neural Turing Machines \cite{graves2014neural}, Memory Networks \cite{weston2015memory}, and DNCs \cite{graves2016differentiable} demonstrate that differentiable memory enables structured reasoning; RAG \cite{NEURIPS2020_6b493230} shows that non-parametric knowledge improves factual accuracy (44.5\% vs 41.5\% on Natural Questions) and supports index updates without retraining; FAISS \cite{8733051} and Sentence-BERT \cite{reimers-gurevych-2019-sentence} make similarity search cheap enough to run in real time. CAEM extends this tradition in three specific ways: the memory is populated only through a four-outcome storage decision governed by a fixed-threshold rule on the calibrated probability composite, with a bounded deferred-reconsider buffer (capacity $K_{\text{def}}$, time-to-live $\tau_{\text{ttl}}$) for borderline items rather than from every retrieval; routing combines similarity with stored-verifier confidence rather than using similarity alone; and an independent safety override forces Tier~3 grounding whenever pre-routing confidence falls below the registered safety floor, regardless of the best memory match.

Continual-learning and self-improvement research supplies the theoretical backbone for iterative refinement. EWC \cite{doi:10.1073/pnas.1611835114} maintains approximately 90\% performance across sequential tasks and supplies the conceptual motivation for bounding cycle-to-cycle drift in adapter space; the explicit $L_{2}$ anchor that approximates EWC under a uniform-Fisher assumption is registered for the inactive full-fine-tune fallback only (Fisher-matrix computation at every cycle boundary on a three-billion-parameter Qwen-2.5-3B-Instruct backbone is prohibitive on a single consumer-grade GPU), while on the shipped low-rank-adapter path the bounded adapter parameter envelope on a frozen backbone acts as the implicit drift bound without the hand-tuned anchor coefficient that the full-parameter fine-tune would require. The parameter-efficient-fine-tuning literature reviewed at \Cref{sec:lit-self-improvement} establishes that low-rank adaptation \cite{hu2022lora} on all linear layers retains general capability across iterative cycles substantially better than full-parameter fine-tuning at matched training budget \cite{biderman2024lora,dettmers2023qlora}, which is the empirical anchor for CAEM's choice of the LoRA primary path with the $L_2$ anchor placed on the adapter parameters rather than on the full backbone. Direct empirical evidence at pre-trained transformer scale shows that the Fisher-weighting advantage of EWC, SI, and MAS over plain $L_2$ collapses once the backbone is a pre-trained language model: \cite{mehta2023empirical} report minimal or negative improvement at BERT-Large scale, and \cite{wu2022pretrained} reproduce the same pattern across five pre-trained backbones, attributed to the flatter loss landscape that pre-training induces. Within-adapter forgetting under sequential fine-tuning is documented by \cite{wang2023olora}, whose orthogonal-subspace decomposition is registered as a future-work upgrade if the multi-modal retention probe panel reports persistent drift on the trajectory's later cycles. Noisy-label learning \cite{natarajan2013learning} proves that training on pools whose purity is well below 100\% can still improve a model, which justifies the registered stricter training threshold $\tau_{\text{train}} > \tau_{\text{store}}$ that lifts training-pool purity above memory purity, where $\tau_{\text{store}}$ is the registered fixed-threshold cut-point on the calibrated probability scale. STaR \cite{NEURIPS2022_639a9a17} demonstrates the bootstrap loop of generate--filter--fine-tune that CAEM generalises by replacing simple correctness filtering with the nine-signal calibrated-probability composite $\ustored$ (per-signal isotonic regression with a per-benchmark shrinkage prior plus an $L_{2}$-regularised boost layer; the contradiction probability $p_{\text{contra}}$ is excluded from the composite signal set because the MiniCheck-only deployment makes it structurally zero), and by persisting verified items in episodic memory rather than only in the adapter. The self-correction critiques of \cite{huang-etal-2023-large} and \cite{kamoi-etal-2024-llms} establish that intrinsic prompt-based self-correction rarely improves accuracy without external feedback, which is exactly the feedback CAEM supplies through the external verifier. Together, these pieces motivate CAEM's cycle-boundary mechanics: retroactive re-verification of stored items against the refreshed adapter, deferred-buffer reconsideration with a three-way branch (promote, drop, keep), pool reweighting with a temperature softmax and a hard share cap, parameter-efficient fine-tuning on the high-purity training pool through the bounded LoRA adapter envelope, and a multi-modal retention guard with floor $\rho_{\min} = 0.93$ on the worst-probe ratio that aborts the cycle and rolls back the adapter (but not the memory) when general capability erodes.

The evaluation benchmark panel is chosen to exercise these mechanisms on substantially disjoint failure modes. FEVER \cite{thorne-etal-2018-fever} tests the external-grounding family and the three-way mapping from SUPPORTS, REFUTES, or NOT ENOUGH INFO onto \textsc{Store}, \textsc{Discard}, or \textsc{Abstain} respectively. TriviaQA \cite{joshi-etal-2017-triviaqa} tests Tier-3 retrieval-augmented generation under noisy evidence on the open-text factoid surface and contributes a held-out test fold to the multi-modal retention probe panel. CommonsenseQA \cite{talmor-etal-2019-commonsenseqa} tests commonsense multi-step reasoning over five-option multiple-choice items and contributes a second held-out test fold to the retention probe panel. TruthfulQA \cite{lin-etal-2022-truthfulqa} and StrategyQA \cite{10.1162/tacl_a_00370} are held out from the SIL training pool as transfer-only benchmarks, exercising the cross-encoder relevance signal on imitative falsehoods and sample-set agreement on implicit multi-hop reasoning respectively. MMLU \cite{hendrycks2021measuring} is the third probe in the multi-modal retention panel, supplying a broad general-knowledge multiple-choice surface that catches catastrophic forgetting on tasks the cycle was not trained on. The headline hallucination metric is the \emph{Composite Hallucination Metric} (CHM), the equal-weighted mean across an eight-subtype failure-mode taxonomy (confident confabulation, factual fabrication, logical fabrication, off-topic, defensive evasion, template leak, false refusal, over-long), aggregated across benchmarks into a geometric Composite Evaluation Score whose five factors capture accuracy, epistemic quality, retention, calibration, and verifier reliability. Specific quantitative results on this panel are reported in Chapter~5.

Taken together, the literature establishes that hallucination remains unsolved despite substantial investment, that parametric scaling is demonstrably insufficient, and that individual methods achieve only limited discriminative accuracy. CAEM's response is to combine verified episodic memory with a four-outcome storage decision under a fixed-threshold rule on the calibrated probability composite, two-stage confidence (cheap pre-routing score plus nine-signal post-generation verifier), a combined-score router with independent safety override, and a constrained self-improvement loop whose hardest protection is a multi-modal retention-guard abort with adapter rollback but memory preservation. Each piece is directly motivated by a specific limitation surfaced in this chapter.

\section{Chapter signpost}
\label{sec:ch2-signpost}

The preliminaries section and the literature review have established the technical foundations on which the architecture rests and have located CAEM's joint intervention against the verification, retrieval, memory-augmented, and continual-learning prior work whose individual contributions do not close the deployment-safety gap registered in \Cref{ch:introduction}. \Cref{ch:requirements} translates the architectural opportunity surfaced by the literature into a registered specification: seven functional requirements paired with named architectural components, five non-functional requirements with empirical anchors, six production-parity design commitments, and the societal, environmental, ethical, standards, project-management, risk, and economic readings that bound the engineering envelope. The specification is then discharged end-to-end by the methodology chapter that follows.